\documentclass[11pt]{article}

\usepackage[margin=1in]{geometry}
\usepackage{amssymb}
\usepackage{booktabs}
\usepackage{enumitem}
\usepackage{xcolor}
\usepackage{listings}
\usepackage[most]{tcolorbox}
\usepackage{hyperref}

\lstset{
  basicstyle=\ttfamily\small,
  columns=fullflexible,
  keepspaces=true,
  breaklines=true,
  frame=single,
  backgroundcolor=\color{gray!8},
  rulecolor=\color{gray!40},
  showstringspaces=false
}

\newtcolorbox{notebox}[1][Note]{colback=blue!4, colframe=blue!50!black, fonttitle=\bfseries, title={#1}}
\newtcolorbox{warnbox}[1][Watch out]{colback=red!5, colframe=red!60!black, fonttitle=\bfseries, title={#1}}

\title{CodePath AI301 --- Open Source Capstone \\ Assignment 1: Issue Selection \\ \large Step-by-Step Guide and Submission Skeleton}
\author{}
\date{}

\begin{document}
\maketitle

\section*{Deliverables at a glance (25 points)}
\begin{center}
\begin{tabular}{@{}lp{6cm}l@{}}
\toprule
\textbf{Where in \texttt{ai301-coursework}} & \textbf{What} & \textbf{Points} \\
\midrule
\texttt{tools/issue-select/} & \texttt{rubric.md}, \texttt{SKILL.md}, \texttt{scope.md}, \texttt{references/} \texttt{evidence-guide.md} & 6 \\
\texttt{beat-1-sandbox/unit-1/eval-run.txt} & one complete harness run (\texttt{--save-run}) & 3 \\
\texttt{beat-1-sandbox/unit-1/selection.md} & four write-up fields & 7 \\
\texttt{beat-1-sandbox/unit-1/selection.md} & chosen issue, verdict, reflection & 9 \\
\bottomrule
\end{tabular}
\end{center}
Submit the link to the \textbf{whole repo}, not a folder.

%----------------------------------------------------------------
\section{Create the course repo}
\begin{enumerate}[label=\textbf{1.\arabic*}]
  \item Open the course template on GitHub $\rightarrow$ \textbf{Use this template}.
  \item Name it \texttt{ai301-coursework}, set it to \textbf{Public}, create it.
  \item Confirm that \texttt{tools/issue-select/} and \texttt{beat-1-sandbox/unit-1/} exist (or will be created when you upload).
\end{enumerate}

%----------------------------------------------------------------
\section{Install the skill}
On Windows, run these in \textbf{Git Bash}, where \texttt{\textasciitilde} is \texttt{C:\textbackslash Users\textbackslash Luis Mendez}.
\begin{lstlisting}
mkdir -p ~/.claude/skills
cp -r path/to/unit-1-materials/skill ~/.claude/skills/issue-select
ls ~/.claude/skills/issue-select
# expect: SKILL.md  scope.md  rubric.md  references/evidence-guide.md
\end{lstlisting}
Edit \texttt{\textasciitilde/.claude/skills/issue-select/scope.md} and set its repo line to:
\begin{lstlisting}
Repo: codepath/pathreview-ai301-fa26-s1
\end{lstlisting}

\begin{notebox}[{One rubric, two uses}]
The installed \texttt{rubric.md} is the \emph{only} copy you edit. Point the eval harness's
\texttt{--rubric} at it, and live mode reads it too, so the two always agree.
\end{notebox}

%----------------------------------------------------------------
\section{Read before you write}
\begin{enumerate}[label=\textbf{3.\arabic*}]
  \item Read \texttt{SKILL.md} (how the skill runs and what it outputs), \texttt{scope.md} (house rules and the repo),
        and \texttt{references/evidence-guide.md} (what counts as evidence).
  \item Read the empty \texttt{rubric.md} and keep its \emph{exact} section layout. The harness and the skill parse it.
  \item Re-read the 4 calibration issues from the activity. They are not scored, so you can use them to sanity-check your checks.
\end{enumerate}

%----------------------------------------------------------------
\section{Draft the rubric}
The motto is \emph{``evidence sources and thresholds, not adjectives.''} A check like ``the issue is well scoped''
is an adjective. A check that names \emph{where to look} and \emph{what number or condition flips the verdict} is a threshold.

\subsection*{Pattern for one check}
Adapt this to the template's format. Fill in the brackets yourself.
\begin{lstlisting}
### Check [N]: [short name]
- Evidence source: [where the skill looks: issue body, labels, comments,
  assignees, linked PRs, files named in the issue, ...]
- Threshold: [the concrete condition, e.g. a count, a presence/absence,
  a time window]
- Verdict effect: [reject | lowers fit]
- Why: [one sentence: what goes wrong for a first-time contributor if
  this check is skipped]
\end{lstlisting}

\subsection*{Questions that lead to checks}
Answer each one with an evidence source and a threshold:
\begin{itemize}
  \item Is someone already working on it? Where would you see that, and how recent must it be to count?
  \item Is the problem reproducible from what is written? What must the issue contain?
  \item Is it small enough for a first issue? What signals size (files touched, subsystems, open questions)?
  \item Is it in scope for Path Review's house rules (see \texttt{scope.md})?
  \item Does it need access, secrets, or maintainer decisions you cannot get?
\end{itemize}

%----------------------------------------------------------------
\section{Evaluate with the cheap loop}
From the \texttt{eval/} directory of the Unit 1 materials (use \texttt{python} or \texttt{py} if \texttt{python3} is missing on Windows):
\begin{lstlisting}
# full run, about $4
python3 run_eval.py --rubric ~/.claude/skills/issue-select/rubric.md

# re-grade only the disagreements, about $0.20 per issue
python3 run_eval.py --rubric ~/.claude/skills/issue-select/rubric.md \
  --only issue-07,issue-12
\end{lstlisting}

\begin{enumerate}[label=\textbf{5.\arabic*}]
  \item After each full run, read the \textbf{per-category tallies}. You must match at least one verdict in \emph{every} category.
  \item For each disagreement, decide: is it one of the ``clear'' 16 (fix the rubric) or an arguable scope call (maybe leave it)?
  \item Change \emph{one} check at a time, then use \texttt{--only} on the issues it should affect.
  \item \textbf{Log every run as you go.} Your ``Run history'' field is graded on order, so keep this table:
\end{enumerate}

\begin{center}
\small
\begin{tabular}{@{}llllp{5.5cm}@{}}
\toprule
\textbf{\#} & \textbf{Type} & \textbf{Score} & \textbf{Categories OK?} & \textbf{What I changed before this run} \\
\midrule
1 & full & \_\_/20 & & initial rubric \\
2 & only & & & \\
3 & & & & \\
4 & & & & \\
\bottomrule
\end{tabular}
\end{center}

\begin{notebox}[The bar]
Aim for \textbf{18/20} with the category floor met. A first run that already scores 20/20 is fine and loses nothing.
\end{notebox}

%----------------------------------------------------------------
\section{Record the final run}
Once the rubric is final, do one \textbf{full} run with \texttt{--save-run}:
\begin{lstlisting}
python3 run_eval.py --rubric ~/.claude/skills/issue-select/rubric.md \
  --save-run eval-run.txt
\end{lstlisting}
\begin{warnbox}
Never hand-edit \texttt{eval-run.txt}. Its header fingerprints your files, so after this run do \textbf{not}
change \texttt{rubric.md}, or the file you upload will no longer match the run. \texttt{--only} and
\texttt{--limit} runs refuse to save.
\end{warnbox}

%----------------------------------------------------------------
\section{Choose your issue (live mode)}
\begin{enumerate}[label=\textbf{7.\arabic*}]
  \item Browse the open issues in \texttt{codepath/pathreview-ai301-fa26-s1} and pick 2--3 candidates.
  \item Run the skill from any directory:
\begin{lstlisting}
claude "issue-select: grade these candidate first issues: <URL1> <URL2> <URL3>"
\end{lstlisting}
  \item Check the output: a ranked read-out (accepted issues with a fit reason, then rejected ones with the check that sank them)
        and a fenced JSON block with \texttt{accept}/\texttt{reject}. \textbf{No JSON block means the skill did not run.}
  \item Of the accepted issues, choose the one that fits you best. Save the verdict output.
  \item \textbf{Do not comment on or claim the issue.} That happens in Unit 2.
\end{enumerate}

%----------------------------------------------------------------
\section{Upload and submit}
\begin{enumerate}[label=\textbf{8.\arabic*}]
  \item In \texttt{tools/issue-select/}: \textbf{Add file $\rightarrow$ Upload files}, then upload the \emph{current} installed
        \texttt{SKILL.md}, \texttt{scope.md}, \texttt{rubric.md}, and \texttt{evidence-guide.md} (inside a \texttt{references/} folder).
  \item In \texttt{beat-1-sandbox/unit-1/}: upload \texttt{eval-run.txt} and the filled \texttt{selection.md}.
  \item Open the repo in a private browser window to confirm it is public, then submit the \textbf{repo} link in the portal.
\end{enumerate}

%----------------------------------------------------------------
\newpage
\section*{Skeleton for \texttt{selection.md}}
\emph{Use the headings that are already in the template file if they differ. Everything in \texttt{[\ldots]} is yours to write, in your own words.}

\begin{lstlisting}
## Chosen issue
Link: [https://github.com/codepath/pathreview-ai301-fa26-s1/issues/N]

Skill verdict:
[paste the ranked read-out line and the JSON entry for this issue]

## Reflection prompts
[answer each prompt from the template here]

## Run history
[In order: run 1 (full, score, category tallies) -> what you changed
 -> run 2 (--only issue-xx, result) -> ... -> final saved run and score.]

## Issue analysis
Issue: [issue-NN]
My rubric's verdict: [accept/reject]   Gold label: [accept/reject]
[Why the rubric read it that way: which check fired or didn't fire, and
 what evidence in the issue it looked at.]

## Check rationale
Quoted check (current wording, identical to the uploaded rubric.md):
> [paste exactly]
[Why this evidence source and this threshold, and what you would lose
 with a looser or stricter threshold.]

## Trade-offs
[What your rubric deliberately accepts or rejects on the arguable
 calls, and the cost of that choice.]
\end{lstlisting}

\subsection*{Final checklist}
\begin{itemize}[label=$\square$]
  \item Repo \texttt{ai301-coursework} is public and made from the template
  \item Skill installed; \texttt{scope.md} Repo line set
  \item $\geq$ 18/20 and category floor met (or the shortfall explained in the write-up)
  \item \texttt{eval-run.txt} from a full \texttt{--save-run}, unedited, and the rubric unchanged since
  \item Skill files uploaded to \texttt{tools/issue-select/}
  \item \texttt{selection.md}: issue link, verdict, reflection, and all four write-up fields
  \item No comment posted on the chosen issue
  \item Whole-repo link submitted
\end{itemize}

\end{document}
